% Volume VI: Adversaries, Platforms, and Automated Judgment
% Part XII. Platforms and Automated Publicity
% Chapter 67: Viral Evidence
% Spine: proof-spines/ch067-spine.md

\chapter{Viral Evidence}
\label{ch:ch067}


Viral media generates attention faster than it generates verification. Let \(V(t)\) denote visibility and \(Q(t)\) verification quality. In the decisive early interval, the characteristic inequality is
\[
\frac{dV}{dt}\gg\frac{dQ}{dt}.
\]
\ToyModel\ The early-interval rate inequality models circulation outrunning verification.
Visibility races upward while provenance, context, and corroboration mature much more slowly.

The consequence is that public judgment often reaches saturation before verification is complete. What millions have already seen, shared, and morally sorted acquires the force of settled evidence even when the record remains epistemically immature. The chapter therefore treats \textbf{viral evidence} as a timing problem: the issue is not that publicity is worthless, but that the rate of circulation can vastly exceed the rate at which evidentiary quality becomes fit for decisive judgment.
